\section{余切复形}\label{src:C.3}
\begin{epigraph}
计算机毫无用处. 它们只能给你答案.
\epigraphauthor{Pablo Picasso}
\end{epigraph}
本章概述 \cite{Ill71} 所引入的概形态射的余切复形的性质, 这一构造将 André \cite{And67} 和 Quillen \cite{Qui68,Qui70} 关于环同态的余切复形的工作整体化. 余切复形的主要优点是它使我们能够描述奇异概形的形变与障碍 (定理 \bookref{C.3.6}).
\subsection{余切复形的性质}\label{src:C.3.1-subsection}
\begin{theorem}\label{src:C.3.1}
对每个概形态射 (相应地, Noether 概形的有限型态射) $f:X\to Y$, 存在由平坦 $\mathcal O_X$-模组成的复形
\[
\mathrm L_{X/Y}:\cdots\to L_{X/Y}^{-1}\to L_{X/Y}^0\to0,
\]
其上同调拟凝聚 (相应地, 凝聚), 且它在 $D_{\mathrm{QCoh}}^-(\mathcal O_X)$ (相应地, $D_{\mathrm{Coh}}^-(\mathcal O_X)$) 中的像也记为 $\mathrm L_{X/Y}$. 它满足以下性质.
\begin{enumerate}[(1)]
\item $\mathrm H^0(X,\mathrm L_{X/Y})\simeq\Omega_{X/Y}$.
\item 态射 $f$ 光滑, 当且仅当 $f$ 局部有限表现且 $\mathrm L_{X/Y}$ 为集中于 $0$ 次的完美复形. 此时有拟同构 $\mathrm L_{X/Y}\xrightarrow{\sim}\Omega_{X/Y}$, 其中 $\Omega_{X/Y}$ 位于 $0$ 次.
\item 若 $f$ 平坦且有限表现, 则 $f$ 为 syntomic 态射 (定义 \bookref{A.3.17}), 当且仅当 $\mathrm L_{X/Y}$ 为集中于 $[-1,0]$ 次的完美复形. 具体地, 若 $f$ 分解为由理想层 $I$ 定义的完全交 $X\hookrightarrow\widetilde X$ 及光滑态射 $\widetilde X\to Y$, 则 $\mathrm L_{X/Y}$ 与 $0\to I/I^2\to\Omega_{\widetilde X/Y}|_X\to0$ 拟同构, 其中 $\Omega_{X/Y}$ 位于 $0$ 次. 若 $f$ 还泛光滑, 则 $L_{X/Y}\simeq\Omega_{X/Y}$.
\item 若
\[
\begin{tikzcd}
X'\arrow[r,"g'"]\arrow[d]&X\arrow[d,"f"]\\
Y'\arrow[r,"g"']\arrow[ur,phantom,"\square",very near start]&Y
\end{tikzcd}
\]
为 Cartesian 图, 且 $f$ 或 $g$ 平坦 (或更一般地, $f$ 与 $g$ 为 Tor 独立的), 则有拟同构 $g'^*\mathrm L_{X/Y}\to\mathrm L_{X'/Y'}$. (注意无须任何平坦性条件, 都有 $g'^*\Omega_{X/Y}\simeq\Omega_{X'/Y'}$.)
\item 若 $X\xrightarrow{f}Y\to Z$ 为概形态射的复合, 则在 $D_{\mathrm{QCoh}}^-(\mathcal O_X)$ 中有正合三角
\[
f^*\mathrm L_{Y/Z}\to\mathrm L_{X/Z}\to\mathrm L_{X/Y}\to f^*\mathrm L_{Y/Z}[1].
\]
它诱导上同调长正合列
\begingroup\footnotesize
\[
\begin{tikzcd}[column sep=8pt,row sep=large]
&\cdots\arrow[r]&\mathrm H^{-2}(\mathrm L_{X/Y})\arrow[dll,out=0,in=180,looseness=1.5]&\\
\mathrm H^{-1}(f^*\mathrm L_{X/Z})\arrow[r]&\mathrm H^{-1}(\mathrm L_{X/Z})\arrow[r]&\mathrm H^{-1}(\mathrm L_{X/Y})\arrow[dll,out=0,in=180,looseness=1.5]&\\
f^*\Omega_{Y/Z}\arrow[r]&\Omega_{X/Z}\arrow[r]&\Omega_{X/Y}\arrow[r]&0,
\end{tikzcd}
\]
\endgroup
延拓通常的微分右正合列 \cite[II.$8.12$]{Har77}. (注意若 $f$ 光滑, 则 $\mathrm H^{-1}(L_{X/Y})=0$, 且 $f^*\Omega_{Y/Z}\to\Omega_{X/Z}$ 单射.)
\end{enumerate}
\end{theorem}
\begin{proof}
概形态射 (更一般地, 赋环 topos 的态射) 的余切复形的定义与构造见 \cite[II.$1.2.3$]{Ill71} 及 \cite[Tag 08T2]{SP}. 对于 (1)--(5), 见 \cite[II.$1.2.4.2$, II.$3.1.2$, II.$3.2.6$, II.$2.2.3$ 及 II.$2.1.2$]{Ill71} 以及 \cite[Tags 08UV, 0D0N, 0FK3, 08QQ 及 08T4]{SP} (注意 \cite[Tag 08RB]{SP} 将朴素余切复形 $\mathrm{NL}_{X/Y}$ 与 $L_{X/Y}$ 联系起来). 对于 (3) 的最后断言, 注意右正合列
\begin{equation}\label{src:C.3.2}
I/I^2\xrightarrow{d}\Omega_{\widetilde X/Y}|_X\longrightarrow\Omega_{X/Y}\longrightarrow0
\end{equation}
也左正合: 因为 $X\hookrightarrow\widetilde X$ 为完全交, $I/I^2$ 为向量丛, 故 $\ker d$ 无挠; 另一方面, 因为 $X\to Y$ 泛光滑, $d$ 泛单射, 所以 $\ker d$ 为挠的. 因此 $\ker d=0$, 且 $L_{X/Y}=[I/I^2\xrightarrow{d}\Omega_{\widetilde X/Y}|_X]$ 与 $\Omega_{X/Y}$ 拟同构.
\end{proof}
\subsection{余切复形的截断}\label{src:C.3.2-subsection}
余切复形的定义依赖单纯方法, 这里不作阐述. 不过, 我们将显式描述它的截断, 这常常已足以应用. 首先, 若 $X\to Y$ 分解为由理想层 $I$ 定义的闭浸入 $X\hookrightarrow P$ 及光滑态射 $P\to Y$, 则 $\mathrm L_{X/Y}$ 在 $[-1,0]$ 次的截断 $\tau_{\geq-1}(\mathrm L_{X/Y})$ 与 $0\to I/I^2\to\Omega_{X/Y}\to0$ 拟同构 (其中 $\Omega_{X/Y}$ 位于 $0$ 次). 若 $X\to Y$ 为 syntomic 态射 (例如光滑态射), 则 $X\hookrightarrow\widetilde Y$ 为正则浸入, $I/I^2$ 为向量丛, 且 $\mathrm L_{X/Y}\simeq\tau_{\geq-1}(\mathrm L_{X/Y})$ (定理 \bookref{C.3.1}(3)).

对仿射概形的态射 $X=\Spec A\to\Spec B=Y$, Lichtenbaum--Schlessinger \cite{LS67} 显式描述 $L_{X/Y}=L_{B/A}$ 的截断 $\tau_{\geq-2}(L_{A/B})$. 选取多项式环 $P=B[x_i]$ (生成元可以有无限多个), 以及核为 $I$ 的 $B$-代数满射 $P\twoheadrightarrow A$. 选取自由 $P$-模 $F$ 及 $P$-模满射 $p:F\twoheadrightarrow I$, 核为 $K=\ker p$. 设 $K'\subseteq K$ 为由 $x,y\in F$ 的 $p(x)y-p(y)x$ 生成的子模. 则 $\tau_{\geq-2}(L_{B/A})$ 与 $A$-模复形
\begin{equation}\label{src:C.3.3}
K/K'\to F\otimes_PA\to\Omega_{P/B}\otimes_PA
\end{equation}
拟同构, 其中最后一项位于 $0$ 次. 另见 \cite[Tag 09CG]{SP}.

对 $i=0,1,2$, 在 $A$-模范畴上定义 $T^i$ 函子:
\[
T^i(A/B,-):=\mathrm H^i(\operatorname{Hom}_A(\mathrm L_{A/B},-))=\mathrm H^i(\operatorname{Hom}_A(\tau_{\geq2}\mathrm L_{A/B},-)),
\]
它们描述概形的形变. 另见 \cite[\S $2.3$]{LS67} 及 \cite[\S $1.3$]{Har10}.
\subsection{代数与概形的扩张}\label{src:C.3.3-subsection}
\begin{definition}\label{src:C.3.4}
概形态射 $X\to S$ 的以拟凝聚 $\mathcal O_X$-模 $J$ 为核的\textit{扩张} 是短正合列
\[
0\to J\to\mathcal O_{X'}\to\mathcal O_X\to0,
\]
其中 $X\hookrightarrow X'$ 为由理想层 $J\subseteq\mathcal O_{X'}$ 定义的概形闭浸入, 满足 $J^2=0$. (注意条件 $J^2=0$ 蕴含 $J\subset\mathcal O_{X'}$ 自然为 $\mathcal O_X$-模.) 平凡扩张为 $X[J]:=(X,\mathcal O_X\oplus J)$, 其环结构由 $J^2=0$ 定义.

扩张的态射是短正合列的态射, 在 $J$ 和 $\mathcal O_X$ 上为恒等映射. 用 $\underline{\operatorname{Exal}}_S(X,J)$ 表示 $X\to S$ 的以 $J$ 为核的扩张范畴, 用 $\operatorname{Exal}_S(X,J)$ 表示同构类集合. 若 $X=\Spec A$ 和 $S=\Spec R$ 仿射, 则写成 $\operatorname{Exal}_R(A,J)$.
\end{definition}
几何上, 扩张是概形的交换图
\[
\begin{tikzcd}
X\arrow[r,hook]\arrow[d]&X'\arrow[dl]\\
\Spec R&
\end{tikzcd}
\]
使得 $J\simeq\ker(\mathcal O_{X'}\to\mathcal O_X)$ 且 $J^2=0$. 扩张集 $\operatorname{Exal}_S(X,J)$ 对 $\mathcal O_X$-模映射 $J\to J'$ 及 $S$-概形的态射 $X'\to X$ 是函子性的, 并继承 $\Gamma(X,\mathcal O_X)$ 上的模结构. 实际上, 群胚 $\underline{\operatorname{Exal}}_S(X,J)$ 为\newterm{Picard 范畴}{Picard category}, 而 $\Sch/S$ 上在 $f:T\to S$ 上的纤维范畴为 $\underline{\operatorname{Exal}}_T(X_T,f^*J)$ 的预叠为\newterm{Picard 叠}{Picard stack}; 见 \cite[III.$1.1.5$]{Ill71} 及 \cite[XVIII.$1.4$]{SGA4}. 给定 $\mathcal O_X$-模正合列 $0\to J'\to J\to J''\to0$, 有正合列
\begingroup\footnotesize
\[
\begin{tikzcd}[column sep=8pt,row sep=large]
0\arrow[r]&\operatorname{Der}_{\mathcal O_S}(\mathcal O_X,J')\arrow[r]&\operatorname{Der}_{\mathcal O_S}(\mathcal O_X,J)\arrow[r]&\operatorname{Der}_{\mathcal O_S}(\mathcal O_X,J'')\arrow[dll,out=0,in=180,looseness=1.5]\\
&\operatorname{Exal}_S(X,J)\arrow[r]&\operatorname{Exal}_S(X,J)\arrow[r]&\operatorname{Exal}_S(X,J'').
\end{tikzcd}
\]
\endgroup
给定 $S$-概形的态射 $f:X\to Y$, 有正合列
\begingroup\footnotesize
\[
\begin{tikzcd}[column sep=8pt,row sep=large]
0\arrow[r]&\operatorname{Der}_{\mathcal O_Y}(\mathcal O_X,J)\arrow[r]&\operatorname{Der}_{\mathcal O_S}(\mathcal O_X,J)\arrow[r]&\operatorname{Der}_{\mathcal O_S}(\mathcal O_Y,f_*J)\arrow[dll,out=0,in=180,looseness=1.5]\\
&\operatorname{Exal}_Y(X,J)\arrow[r]&\operatorname{Exal}_S(X,J)\arrow[r]&\operatorname{Exal}_S(Y,f_*J).
\end{tikzcd}
\]
\endgroup
见 \cite[$0.20.2.3$]{EGA} 和 \cite[III.$1.2.4.3$, III.$1.2.5.4$]{Ill71}. 上述第二个图的第一行, 可由对右正合列 $f^*\Omega_{Y/S}\to\Omega_{X/S}\to\Omega_{X/Y}\to0$ 应用 $\operatorname{Hom}_{\mathcal O_X}(-,J)$, 并使用以下等同实现:
\[
\begin{aligned}
\operatorname{Hom}_{\mathcal O_X}(\Omega_{X/Y},J)&=\operatorname{Der}_{\mathcal O_Y}(\mathcal O_X,J),\\
\operatorname{Hom}_{\mathcal O_X}(\Omega_{X/S},J)&=\operatorname{Der}_{\mathcal O_S}(\mathcal O_X,J),\\
\operatorname{Hom}_{\mathcal O_X}(f^*\Omega_{Y/S},J)&=\operatorname{Hom}_{\mathcal O_Y}(\Omega_{Y/S},f_*J)=\operatorname{Der}_{\mathcal O_S}(\mathcal O_Y,f_*J).
\end{aligned}
\]
\subsection{余切复形与形变理论}\label{src:C.3.4-subsection}
\begin{theorem}\label{src:C.3.5}
若 $X\to Y$ 为概形态射, $J$ 为拟凝聚 $\mathcal O_Y$-模, 则有自然同构
\[
\operatorname{Exal}_Y(X,J)\simeq\operatorname{Ext}_{\mathcal O_X}^1(\mathrm L_{X/Y},J).
\]
\end{theorem}
\begin{proof}
见 \cite[III.$1.2.3$]{Ill71}.
\end{proof}
对正合列 $0\to J'\to J\to J''$ 应用 $\operatorname{Hom}_{\mathcal O_X}(\mathrm L_{X/Y},-)$, 并对正合三角 $f^*\mathrm L_{Y/Z}\to\mathrm L_{X/Z}\to\mathrm L_{X/Y}$ 应用 $\operatorname{Hom}_{\mathcal O_X}(-,J)$, 可将上述两个六项正合列延拓为长正合列. 当 $X=\Spec A\to\Spec B=Y$ 为仿射概形的态射时, 使用 \S\bookref{C.3.2-subsection} 的 $T^i$ 函子, 上述等同变为 $\operatorname{Exal}_B(A,J)=T^1(A/B,J)$, 它也可由截断余切复形 \bookeqref{C.3.3} 显式给出; 见 \cite[$4.2.2$]{LS67} 及 \cite[定理 $5.1$]{Har10}. $T^i$ 函子如何将上述六项序列延拓为九项序列, 见 \cite[$2.3.5$--$6$]{LS67} 及 \cite[定理 $3.4$--$5$]{Har10}.
\begin{theorem}\label{src:C.3.6}
考虑形变问题
\[
\begin{tikzcd}
X\arrow[r,hook,dashed]\arrow[d,"f"']&X'\arrow[d,dashed,"f'"]\\
Y\arrow[r,hook,"i"']&Y',
\end{tikzcd}
\]
其中 $f:X\to Y$ 为概形态射, $i:Y\hookrightarrow Y'$ 为由理想层 $J\subseteq\mathcal O_{Y'}$ 定义的概形闭浸入, 满足 $J^2=0$. 形变是使上述图为 Cartesian 图的态射 $f':X'\to Y'$, 而形变态射是 $Y'$ 上的态射, 在 $X$ 上限制为恒等态射.
\begin{enumerate}[(1)]
\item 形变 $f':X'\to Y'$ 的自同构群同构于 $\operatorname{Ext}_{\mathcal O_X}^0(\mathrm L_{X/Y},f^*J)$.
\item 若存在形变, 则形变集为 $\operatorname{Ext}_{\mathcal O_X}^1(\mathrm L_{X/Y},f^*J)$ 下的 torsor.
\item 存在元素 $\operatorname{ob}_X\in\operatorname{Ext}_{\mathcal O_X}^2(L_{X/Y},f^*J)$, 使得存在形变, 当且仅当 $\operatorname{ob}_X=0$.
\end{enumerate}
\end{theorem}
\begin{proof}
见 \cite[III.$2.1.7$]{Ill71} 及 \cite[Tag 08UZ]{SP}. 关于使用截断余切复形的仿射情形描述, 另见 \cite[$4.2.5$]{LS67} 及 \cite[定理 $10.1$]{Har10}.
\end{proof}
作为检验, 若 $f:X\to\Spec A$ 为 syntomic 态射 (例如光滑态射), $A'\twoheadrightarrow A$ 为核 $J$ 的平方为零的环满射, 则等同
\[
\operatorname{Ext}_{\mathcal O_X}^i(\mathrm L_{X/A},f^*J)=\operatorname{Ext}_{\mathcal O_X}^i(\Omega_{X/A},J)
\]
重新给出命题 \bookref{C.2.4}.
\begin{remark}\label{src:C.3.7}
其他形变问题也有类似结果. 例如, 对形变问题
\[
\begin{tikzcd}[column sep=small]
X\arrow[r,hook]\arrow[dr,"f"']\arrow[ddr]&X'\arrow[dr,dashed]\arrow[ddr]&\\
&Y\arrow[r,hook]\arrow[d]&Y'\arrow[d]\\
&Z\arrow[r,hook]&Z',
\end{tikzcd}
\]
其中水平态射为分别由平方为零的理想层 $J_X$, $J_Y$ 和 $J_Z$ 定义的闭浸入, 则自同构, 形变和障碍由 $\operatorname{Ext}_{\mathcal O_X}^i(Lf^*\mathrm L_{Y/Z},J_X)$ 分类, 其中 $i=-1,0,1$ \cite[III.$2.2.4$]{Ill71}.
\end{remark}
